\subsection{位数6のOre円錐と曲率枝の排除}
\label{subsec:cubic-sixth-root-spectrum}
\begin{lemma}[有限重みから強制されるOre円錐]
\label{lem:cubic-finite-spectrum-forced-cone}
三固有値候補の単純固有基底を $v$ とする。
各他固有空間 $V_j$ に対し、重み $\lambda_v\lambda_j$ の二次tensor空間が
$vV_j+V_jv$ だけで、その関係次元が $\dim V_j$ であるとする。
また重み $\lambda_v^2$ に関係がなく、三次重み1の全指数組が $v$ の個数0であるとする。
この候補はH3を与えない。
\end{lemma}
\begin{proof}
整域性から各mixed関係の両projectionは同型である。
従って $v$ は残る三生成元を可逆一次変換 $\sigma$ で正規化する。
$v^2$ を含む関係はなく、残りの三本の二次関係は他の三生成元だけからなる。
その三生成元のquadratic提示を $A$ とすると、
$v r=\sigma(r)v=0$ と整域性により $\sigma$ は三次元の関係空間を保存する。
生成元上可逆なので $\sigma\in\operatorname{Aut}_{\rm gr}A$、全二次提示から $B=A[v;\sigma]$ を得る。
三次重み条件により $G\in A_3$。
$\theta=\mu_B^{-1}$ は $A$ と $kv$ を保存し、$\theta(v)=\lambda_v^{-1}v$。
補題\ref{lem:cubic-skew-cone-vertex}をscalar normalizer版で適用でき、
$\pi k[v]$ の無限高次自己ExtからH3でない。
\end{proof}

\begin{proposition}[位数6の四つの円錐スペクトルの排除]
\label{prop:cubic-sixth-root-four-cones}
有限リストの次の四指数組はH3を持たない。
\[
 (N;a,b,c)=(6;1,3,4),(6;2,3,5),(6;4,1,3),(6;5,2,3).
\]
\end{proposition}
\begin{proof}
最初の二つでは $\lambda_b=-1$、後二つでは $\lambda_c=-1$ の単純基底を $v$ と取る。
全三次共鳴は $v$ を含まず、$v^2$ の関係はなく、
$vU+Uv$ の関係数2、もう一つの単純基底とのmixed関係数1が、それぞれ孤立した重み空間にある。
公開検証器で全四組の指数と関係重みを確認し、前補題を適用する。
$v$ の正常化は $-1$ 倍なので、可換中心coneだけに限定した補題では足りず、
先のscalar normalizer版を用いている。
\end{proof}

\begin{theorem}[位数6の二つの曲率スペクトルの排除]
\label{thm:cubic-sixth-root-curvature-empty}
有限リストの $(6;1,0,4)$、$(6;5,0,2)$ もH3を持たない。
\end{theorem}
\begin{proof}
原始6乗根 $\xi$ により $\mu_{B,s}=(\xi,\xi,1,\xi^{-2})$ と書く。
二次関係の重みは $Ub+bU$ に2、$Uc+cU$ に2、$bc,cb$ に1、$U^2+kc^2$ に1である。
従って
\[
 bu=f(u)b,\quad cu=g(u)c,\quad bc=qcb,\quad r(U)=\lambda c^2
\]
と書ける。$f,g$ は可逆、$q\ne0$。
Frobenius middleで $U^2+kc^2$ 対 $bc,cb$ のrank1 blockを巡回すると、
二つのmixed graphの可逆小blockを得るので、$r$ の二元係数行列は非退化である。
従って $T=T(U)/(r)$ は形式的に二元AS2平面を定める。

$\lambda=0$ なら、三次overlapと $\dim B_3=20$ により
$fg=gf$、$f,g\in\operatorname{Aut}T$ となり、$B$ は二元AS2上の二つの自己同型型Ore拡大。
可換非半単純作用は命題\ref{prop:cubic-commuting-Ore-no-shear}で排除、
可換半単純作用と対角量子作用はskewとなり、三固有値H3を持たない。
Jordan基底は命題\ref{prop:cubic-Jordan-Ore-resonance}の必要固有値1条件が
重複中山固有値 $\xi\ne1$ と矛盾する。
反交換基底の反対角作用は、位数5の証明の $d=0$ と同じ中心固定・奇次数係数比較で正常核が零。
よってこの枝は空である。

$\lambda\ne0$ とする。
$b,c$ による共役は関係 $r=\lambda c^2$ を保つため、純 $U^2$ のtensorとして
\[
 f(r)=q^2r,\qquad g(r)=r.
\]
従って $\operatorname{hdet}_T(f)=q^2$。
$b$ を除く三生成元が生成する $A$ に対し $B=A[b;\sigma]$、
$\sigma|_U=f$、$\sigma(c)=q c$ となる。
$A$ は $B/(b)$ と同じquadratic提示を持ち、正規正則一次元による商なのでAS Gorenstein次元3。
また $c$ は $A$ 上正規正則で $A/(c)=T$。
この正規一次元のequivariant Rees lemmaで最高Extのcharacterを比較すると
\[
 \operatorname{hdet}_A(\sigma)=q\operatorname{hdet}_T(f)=q^3.
\]
Ore中山公式と $\mu_B(b)=b$ より $q^3=1$ を得る。

三次共鳴は $b^3,U^2c,c^3$ だけ。
$c$ は $B$ 全体で正規なので、整域商のためには $G$ の $b^3$ 係数が非零でなければならない。
正常式 $Gc=\xi^2cG$ の $b^3$ 成分は $q^3=\xi^2$ を要求する。
$\xi^2\ne1$ であり、先の $q^3=1$ と矛盾する。
\end{proof}

\begin{proposition}[位数6の反対重複固有値の正常核消滅]
\label{prop:cubic-sixth-root-opposite-repeat-empty}
指数組 $(6;3,2,4)$ の候補は $W_1(B)=0$ である。
\end{proposition}
\begin{proof}
固有値を $(-1,-1,\omega,\omega^2)$ と書く。
関係空間は $U^2+kbc+kcb$ に2、$Ub+bU$ に2、$Uc+cU$ に2で、
$b^2,c^2$ を含まない。
従って $\deg U=0,\deg b=1,\deg c=-1$ という補助 $\mathbb Z$ 次数が成立する。
不変三次は $G=C b^3+D c^3$ だけである。
正常式 $Gb=\omega^{-1}bG$ の補助次数4成分から
$(1-\omega^{-1})C b^4=0$、整域性から $C=0$。
同様に $Gc=\omega^{-2}cG$ の補助次数 $-4$ 成分から $D=0$。
中山の $U$ 上の単冪部分はこの比較を変えない。
\end{proof}

\begin{theorem}[位数6の二AS2平面型の排除]
\label{thm:cubic-sixth-root-two-AS2-empty}
指数組 $(6;2,1,4)$、$(6;4,2,5)$ の候補にもH3は存在しない。
\end{theorem}
\begin{proof}
前者で $\mu_{B,s}|_U=\omega I$、$\omega$ を原始3乗根とする。
関係重みから、純 $U$ 関係 $r$ が一本、
$b^2=\kappa c^2$ ($\kappa\ne0$) が一本、
$bu=f(u)b,cu=g(u)c$ ($u\in U$) が各二本となる。
$r$ は非退化で $T=T(U)/(r)$ はAS2、
$V=k\langle b,c\rangle/(b^2-\kappa c^2)$ もAS2で、線型基底変更により反交換平面である。
全関係は $U$ 次数と $b,c$ 次数の二重次数を保つ。
不変三次は $G=H+K$、$H\in T_3,K\in V_3$ と書ける。
一方が零なら純AS2三次の因子障害で整域商にならないので、両方非零とする。

二重次数を使って $U$ との正常式を比較すると
$Hu=\mu_B^{-1}(u)H$ を得る。
可換 $T$ ではこの式が $\mu_B|_U=1$ を強制し、$\omega\ne1$ と矛盾。
Jordanでは正規一次元 $x$ に対する多項式の平行移動式の最高係数から、
normalizerの $x$ 固有値は1でなければならず、同じく矛盾する。
従って $T$ は量子平面 $xy=p yx$ である。
その自己同型は $p=-1$ の反対角を除き対角。
反対角は二つの固有値が相異なるため、現在の重複中山にはなれない。
よって実際に $\mu_B|_U=\omega I$ である。
$H=\sum_i H_i x^iy^{3-i}$ の正常式は、非零項ごとに
\[
 p^i=\omega^{-1},\qquad p^{-(3-i)}=\omega^{-1}
\]
を与える。従って $p^3=1$、$p\ne1$。
$f,g$ はこの $T$ の自己同型なので対角であり、
$bx=\alpha xb,by=\beta yb,cx=\gamma xc,cy=\delta yc$ と書ける。
$b^2=\kappa c^2$ と整域性から $\alpha^2=\gamma^2,\beta^2=\delta^2$。

従って $x,y$ は $V$ をそれぞれ対角自己同型 $F,J$ で正規化し、
$F(b)=\alpha^{-1}b,F(c)=\gamma^{-1}c$、
$F^2=\alpha^{-2}I$、$\operatorname{hdet}_V F=\alpha^{-2}$。
同じ全二次提示から
$B=V[x;F][y;\widetilde J]$、$\widetilde J(x)=p^{-1}x$ である。
Ore中山公式から
\[
 \omega=\mu_B(x)/x=p\alpha^{-2},\qquad
 \alpha^6=(p/\omega)^3=1.
\]
一方、$K$ に対する $x$ の正常式は $F(K)=\omega K$。
次数3で $F^2$ は $\alpha^{-6}$ 倍なので、$K\ne0$ から
$\alpha^{-6}=\omega^2$ を得る。これは $\alpha^6=1$ と矛盾する。
後者の指数組も $\omega$ の共役と $b,c$ の交換で同じ議論となる。
\end{proof}

\begin{corollary}[現在残る三固有値の17スペクトル]
\label{cor:cubic-three-eigenvalue-17-list}
以上の排除後、三固有値H3の必要リストは、位数3の3個、位数4の12個、
および位数6の $(6;2,0,5),(6;4,0,1)$ の2個、合計17個に限られる。
Galoisラベルは10軌道である。残る候補の存在やH3十分性は未完成である。
\end{corollary}
